\documentclass[11pt]{article}

\usepackage[margin=1in]{geometry}
\usepackage{booktabs}
\usepackage{array}
\usepackage{hyperref}
\usepackage{xcolor}

\title{CAI6734 Weekly Report 1}
\author{Omer Kahveci}
\date{September 29, 2026}

\begin{document}

\maketitle

\section{Project Objective}
The goal of this project is to train a search agent capable of navigating structured EHR data. This week's scope was 
to reproduce the methodology presented in the LA-CDM paper~\cite{baniharouni2025language}. The goal of this task was to 
improve our understanding of how a pre-trained model can be adapted to support clinical decision-making.

The paper fine-tuned a pre-trained Qwen model using Low-Rank Adaptation (LoRA), a parameter-efficient method that
freezes the original model weights and trains a smaller set of adapter parameters. These adapters were optimized
using supervised fine-tuning and Group Relative Policy Optimization (GRPO) to improve hypothesis generation,
confidence calibration, and diagnostic decision-making.

The model was trained and evaluated on a cohort of 2,400 MIMIC-IV admissions across four acute abdominal conditions: 
appendicitis, cholecystitis, diverticulitis, and pancreatitis. The cohort was divided into 1,920 training cases, 240 
validation cases, and 240 test cases.

\section{Accomplishments}

I obtained credentialed access to MIMIC-IV and used version 2.2 of the dataset on HiPerGator to reproduce the study's 
methodology. I used Meta-Llama-3-8B-Instruct in place of Qwen because Chinese-developed models were not permitted on 
HiPerGator. Since the paper's codebase was open source, I was able to clone it and save a copy on my personal account.

I modified the original LA-CDM codebase to support Llama and the HiPerGator environment. These changes included 
replacing Qwen-specific model paths, resolving dependency conflicts, adding Slurm job scripts, adapting the output
parser to Llama's response format, reducing GPU memory usage, and adding validation scripts for both the cohort and 
prepared datasets. I also updated the evaluation pipeline to report diagnostic accuracy, calibration, test usage,
diagnostic cost, and disease-specific performance.

I extracted a cohort from MIMIC-IV and MIMIC-IV-Note version 2.2 using the same four abdominal conditions studied in 
LA-CDM. The final validated cohort contained 2,400 admissions: 957 appendicitis cases, 648 cholecystitis cases,
257 diverticulitis cases, and 538 pancreatitis cases. I divided the cohort into 1,920 training cases, 240 validation 
cases, and 240 test cases. I then generated Llama-based patient-history summaries and validated the resulting files for
missing summaries, duplicate admissions, malformed fields, diagnosis leakage, and overlap between dataset splits.

Finally, I evaluated the untrained Llama model on all 240 test cases to establish a zero-shot baseline. The model 
achieved 8.75\% overall accuracy and produced no valid final diagnosis for 87.92\% of cases. However, it achieved 
72.41\% accuracy among the cases for which it completed a valid diagnosis. This suggests that the model's main 
zero-shot limitation was following and completing the required interaction protocol rather than identifying the correct 
condition. I also validated the LoRA training pipeline through short training, checkpoint, adapter-loading, and 
production-scale memory tests in preparation for the full training run.

\begin{thebibliography}{1}

\bibitem{baniharouni2025language}
David Bani-Harouni, Chantal Pellegrini, Ege {\"O}zsoy, Nassir Navab, and
Matthias Keicher.
\newblock Language agents for hypothesis-driven clinical decision making with
reinforcement learning.
\newblock In \emph{International Conference on Learning Representations}, 2026.
\newblock \url{https://arxiv.org/abs/2506.13474}.

\end{thebibliography}

\end{document}
